\documentclass[conference]{IEEEtran}

\usepackage[T1]{fontenc}
\usepackage[utf8]{inputenc}
\usepackage[brazilian]{babel}
\usepackage{cite}
\usepackage{amsmath,amssymb,amsfonts}
\usepackage{graphicx}
\usepackage{array}
\usepackage{tikz}
\usetikzlibrary{positioning,calc}

% Permite hifenização em português para evitar grandes espaços no texto justificado.
\hyphenpenalty=250
\exhyphenpenalty=1000

\title{Classificação Incremental e Auditável de Crimes em Notícias Institucionais com Memória WNN e LLM Residual}

\author{
\IEEEauthorblockN{Jefferson S. dos Anjos}
\IEEEauthorblockA{
\textit{Centro Federal de Educação Tecnológica Celso Suckow da Fonseca (CEFET/RJ)}\\
Rio de Janeiro, RJ, Brasil\\
jefferson.anjos.1@aluno.cefet-rj.br}
}

\begin{document}

\maketitle

\begin{abstract}
Bases textuais institucionais crescem continuamente e tornam custosa a classificação temática, a revisão humana e, sobretudo, o envio indiscriminado de tokens a modelos de linguagem. Este trabalho propõe uma metodologia incremental e auditável para classificar o crime canônico principal em notícias públicas da Polícia Federal. A abordagem combina fundação temporal, memória lexical binária inspirada em WNN/WiSARD e revisão residual por LLM. Na rodada concluída, a memória classificou diretamente 53,48\% da reserva, evitando igual proporção de chamadas e restringindo a LLM aos 46,52\% de casos residuais. Contra as tags institucionais mapeadas após a decisão, a WNN obteve precisão de 85,29\%, revocação de 49,91\% e medida-F1 de 62,97\%. Os resultados preservam evidências lexicais, versões da memória e decisões residuais para auditoria.
\end{abstract}

\begin{IEEEkeywords}
classificação textual, aprendizado incremental, redes neurais sem pesos, WiSARD, auditoria, modelo de linguagem
\end{IEEEkeywords}

\section{Introdução}

Instituições públicas que produzem ou analisam informação em larga escala lidam continuamente com grandes volumes de textos. Notícias, registros administrativos, documentos técnicos e comunicados institucionais acumulam evidências relevantes para a análise de políticas públicas, mas nem sempre estão organizados em categorias estáveis. O desafio, portanto, não se limita ao armazenamento dos documentos: consiste em transformar um fluxo textual contínuo em informação temática comparável, auditável e reutilizável.

Esse processo é dificultado pela heterogeneidade da base, pelo seu crescimento ao longo do tempo e pela coexistência de fenômenos recorrentes e emergentes. A variação lexical, as mudanças de enfoque e a presença de atributos acidentais --- como localidades, nomes de operações e entidades --- podem induzir categorias artificiais. Embora a rotulagem humana seja onerosa e pouco escalável, uma classificação automática sem controles pode reproduzir ruídos, transformar metadados em classes e comprometer a estabilidade da taxonomia.

Modelos de linguagem de grande porte podem apoiar a interpretação de textos e a identificação de evidências, mas sua aplicação indiscriminada sobre toda a base aumenta o custo de inferência e reduz a previsibilidade do processo. Este trabalho explora uma divisão de responsabilidades: a memória WNN trata documentos recorrentes para os quais existem evidências lexicais suficientes, enquanto a LLM revisa apenas casos residuais, ambíguos ou semanticamente inéditos. Os aprendizados validados nessa revisão podem ser convertidos em marcadores auditáveis e incorporados de forma controlada à memória.

Assim, o objetivo é responder à seguinte questão: \textit{como organizar e classificar incrementalmente uma base textual institucional em crescimento, reduzindo o custo de inferência por LLM e preservando auditabilidade, adaptação a novos temas e estabilidade da taxonomia?} A metodologia é aplicada a notícias públicas da Polícia Federal e tem como alvo o crime canônico principal de cada documento. Sua contribuição é propor um ciclo reproduzível de fundação temática, classificação incremental, revisão residual e atualização controlada da memória, e não apenas uma classificação pontual da base analisada.

\section{Trabalhos relacionados}

Rangel et al.~\cite{rangel2016ssw} propõem uma WiSARD semissupervisionada para categorização de textos sociais. A abordagem explora a rapidez de treinamento e o aprendizado incremental da arquitetura, empregando uma medida de confiança para selecionar exemplos não rotulados que podem realimentar o classificador. Nos experimentos com conjuntos binários de sentimento, o estudo relata acurácias competitivas e tempo de treinamento inferior aos mecanismos semissupervisionados usados na comparação. Sua contribuição fundamenta, portanto, o uso de redes neurais sem pesos em cenários textuais com poucos rótulos e fluxo de exemplos.

Entretanto, o objetivo de Rangel et al.~\cite{rangel2016ssw} é classificar textos em categorias previamente definidas. O trabalho não tem como foco a descoberta de uma fundação temática, a validação de categorias emergentes nem o controle da evolução de uma taxonomia institucional. Essas necessidades são centrais em notícias públicas: termos como localidades, nomes de operações e órgãos podem ser frequentes sem constituir, por si só, crimes canônicos.

Carneiro~\cite{carneiro2017weightless} investiga resultados teóricos para classificadores sem pesos e sua aplicação à linguística computacional. A tese demonstra que a técnica de \textit{bleaching} aumenta a robustez da WiSARD diante de grandes volumes de dados sem reduzir sua capacidade teórica de generalização e a aplica à etiquetagem gramatical multilíngue. O estudo também evidencia que a representação de entrada é determinante para o desempenho de uma WNN, oferecendo sustentação para a retina textual e para o vocabulário versionado empregados neste trabalho.

Todavia, a aplicação de Carneiro~\cite{carneiro2017weightless} tem como alvo etiquetas linguísticas e a análise das propriedades do classificador. Ela não aborda a construção de uma taxonomia temática de crimes, a governança de documentos raros ou ambíguos, nem o uso seletivo de um modelo de linguagem para interpretar casos sem evidência suficiente na memória.

No campo da supervisão fraca, a família Snorkel usa funções de rotulagem para produzir sinais programáticos e dados de treinamento quando há poucos rótulos manuais~\cite{ratner2017snorkel,ratner2018drybell,ratner2020snorkel}. De modo complementar, Fries et al.~\cite{fries2021ontology} mostram que ontologias, regras e conhecimento de domínio podem apoiar classificações especializadas. A proposta compartilha o emprego de sinais interpretáveis, mas difere porque os discriminadores lexicais não são somente uma etapa intermediária para treinar outro classificador: eles permanecem como memória operacional, versionada e auditável.

Na organização de coleções textuais, LDA~\cite{blei2003lda}, BERTopic~\cite{grootendorst2022bertopic}, HDBSCAN~\cite{mcinnes2017hdbscan} e Sentence-BERT~\cite{reimers2019sbert} oferecem mecanismos para modelar tópicos, agrupar documentos ou estimar proximidade semântica. Neste trabalho, esses recursos apoiam exclusivamente a fundação exploratória. Os clusters não são assumidos como categorias finais: precisam ser consolidados em crimes canônicos e convertidos em discriminadores que preservem a justificativa da decisão.

Estudos sobre LLMs em supervisão fraca e classificação textual mostram seu potencial para formular funções de rotulagem, interpretar textos e aproximar clusterização e classificação~\cite{smith2022lml,guan2023labelfunctions,huang2024clustering}. Em particular, Guan et al.~\cite{guan2023labelfunctions} propõem gerar e filtrar funções de rotulagem com modelos de linguagem. Sosnovikov et al.~\cite{sosnovikov2025lgiwl} empregam a LLM para sintetizar e refinar iterativamente essas funções a partir de erros de um classificador posterior. Essas abordagens aproximam-se da geração de discriminadores interpretáveis, mas as regras servem para supervisionar outro modelo. Nesta metodologia, os discriminadores permanecem como a memória operacional de classificação e a LLM atua somente sobre os casos residuais.

O encaminhamento seletivo por confiança também aparece em arquiteturas de roteamento de LLMs. Shah e Shridhar~\cite{shah2025select} combinam uma taxonomia de tarefas com uma cascata adaptativa que encaminha consultas de baixa confiabilidade para modelos mais capazes. A finalidade é semelhante à redução de custo de inferência, mas o objeto é distinto: o trabalho seleciona LLMs para consultas gerais, enquanto a proposta encaminha notícias institucionais entre uma memória lexical de crimes e uma única etapa residual de interpretação.

No aprendizado incremental, Xia et al.~\cite{xia2021incremental} formalizam a classificação textual com múltiplas rodadas de novas classes e preservação do desempenho nas classes anteriores. Na construção de taxonomias, Tuan et al.~\cite{tuan2016temporal} propõem atualizar relações taxonômicas com evidências temporais, evitando reconstruir toda a estrutura a cada atualização. A construção de taxonomias com LLMs, \textit{prompting} e ajuste fino também é investigada por Balakrishnan~\cite{balakrishnan2025taxonomy}. Esses estudos tratam, respectivamente, da entrada de classes, da atualização da taxonomia ou da sua geração, mas não associam esses processos a uma memória lexical binária, a critérios explícitos de suficiência de evidência e à revisão residual por LLM.

Por sua vez, o princípio de converter evidências observadas em novos marcadores aproxima-se do \textit{bootstrapping} de padrões de Snowball~\cite{agichtein2000snowball}; a diferença é que o objetivo aqui é o aprendizado incremental de classes temáticas, e não a extração de relações.

Assim, a lacuna identificada situa-se entre a eficiência e o aprendizado incremental de WNNs, a geração de regras por LLMs, a atualização de categorias e o roteamento sensível à confiança. A contribuição desta proposta é integrá-los em um ciclo único e auditável: uma fundação formada exclusivamente por documentos históricos consolida crimes canônicos e seus discriminadores; a memória WNN classifica os casos recorrentes; e somente os documentos sem evidência suficiente são encaminhados à LLM. As decisões residuais podem reforçar marcadores existentes, propor candidatos a novos crimes ou registrar exceções, sempre com controle de vocabulário, versão da memória e trilha de auditoria. Portanto, o destaque não é uma nova arquitetura WNN, mas a organização operacional e rastreável desses componentes para uma base institucional em evolução.

\section{Metodologia}

\subsection{Visão geral}

Propõe-se uma metodologia incremental para classificar notícias institucionais pelo crime canônico principal. A base completa é ordenada pela data de publicação e dividida por um corte cronológico em duas partes disjuntas: a fundação contém os 10\% mais antigos dos documentos, e a reserva incremental contém os 90\% posteriores. A fundação é usada para descobrir e consolidar a taxonomia inicial sem acesso a documentos futuros; a reserva é processada em lotes cronológicos de 500 notícias para avaliar a cobertura da memória WNN, os encaminhamentos residuais, a evolução temporal dos crimes e o aprendizado controlado de novos marcadores.

O corpo da notícia é a fonte textual principal para a classificação, pois descreve o fato jurídico e suas evidências. O título pode ser usado como evidência secundária quando o corpo não ativa discriminadores suficientes. As tags institucionais são tratadas apenas como pistas auxiliares: elas podem reforçar uma classe já sustentada por evidência textual, mas não criam isoladamente uma classificação nem uma posição na memória.

A memória adotada é uma implementação lexical binária inspirada em princípios de WNN/WiSARD, e não uma reprodução da WiSARD clássica. Em particular, ela não emprega o mapeamento pseudoaleatório da retina para um conjunto de RAMs endereçáveis. A implementação associa posições lexicais estáveis a discriminadores de crime, constrói vetores binários para os documentos e decide por pontuação de máscaras, cobertura, confiança e margem. A inspiração nas redes neurais sem pesos orienta a representação e a atualização incremental, enquanto as regras de decisão são explicitadas e auditadas no domínio textual.

\subsection{Nomenclatura e representação textual}

Para evitar ambiguidade com a terminologia original da WiSARD, este trabalho utiliza os termos apresentados na Tabela~\ref{tab:nomenclatura-wnn}.

\begin{table}[ht]
\centering
\caption{Nomenclatura da memória WNN aplicada a textos.}
\label{tab:nomenclatura-wnn}
\begin{tabular}{p{0.28\linewidth}p{0.64\linewidth}}
\hline
\textbf{Termo} & \textbf{Definição operacional} \\
\hline
Texto semântico & Representação pré-processada do corpo da notícia, com remoção de ruído e retenção de termos substantivos relevantes. \\
Retina textual & Camada de codificação que projeta o texto semântico em um vetor binário de posições lexicais fixas. O termo é usado por analogia à retina da WiSARD; não se refere a imagem. \\
Vocabulário versionado & Sequência ordenada $V^{(t)}=(a_1,a_2,\ldots,a_{n_t})$ de posições lexicais normalizadas na versão $t$ da memória. \\
Posição lexical $a_i$ & Índice estável associado a uma palavra discriminativa normalizada, por exemplo, $a_{12}=\textit{cartao}$. \\
Vetor binário da notícia & Vetor $\mathbf{x}(d)$ que indica quais posições do vocabulário foram ativadas pelo documento $d$. \\
Máscara de crime & Vetor binário contínuo associado a um crime canônico, indicando as posições lexicais que fundamentam sua identificação. \\
Imagem binária ou cartão perfurado & Representação visual da máscara ou do vetor da notícia. As quebras em linhas servem apenas para caber no espaço da figura; não representam novas dimensões nem novas classes. \\
Discriminador & Conjunto pequeno de posições e termos que, quando ativados em combinação, fornece evidência para uma máscara de crime. \\
\hline
\end{tabular}
\end{table}

Seja $V^{(t)}=(a_1,a_2,\ldots,a_{n_t})$ o vocabulário versionado na iteração $t$. Para uma notícia $d$, a retina textual produz o vetor binário

\begin{equation}
x_i(d)=
\begin{cases}
1, & \text{se o termo associado a } a_i \text{ ocorre em } d;\\
0, & \text{caso contrário.}
\end{cases}
\end{equation}

Logo, $\mathbf{x}(d)\in\{0,1\}^{n_t}$. Cada crime canônico $c$ possui uma máscara binária contínua

\begin{equation}
\mathbf{m}_{c}^{(t)}=(m_{c1},m_{c2},\ldots,m_{cn_t})\in\{0,1\}^{n_t},
\end{equation}

em que $m_{ci}=1$ indica que a posição $a_i$ fundamenta o crime $c$, e $m_{ci}=0$ indica ausência dessa associação. Inspirada na retina binária e nos discriminadores RAM da WiSARD~\cite{rangel2016ssw,carneiro2017weightless}, a memória é apresentada como cartões perfurados: cada cartão visualiza a máscara binária de um crime sobre o mesmo vocabulário versionado. Essa é uma visualização interpretável proposta neste trabalho; não substitui a arquitetura RAM clássica. Uma nova linha exibida na grade representa somente a continuação do vetor na linha seguinte.

Por exemplo, se $a_{12}=\textit{cartao}$, $a_{31}=\textit{perfurado}$ e $a_{44}=\textit{fraude}$, a máscara de um crime financeiro terá valor $1$ nessas posições. A expressão \textit{cartão perfurado} não exige uma posição adicional: ela é representada pela ativação conjunta de $a_{12}$ e $a_{31}$.

A Figura~\ref{fig:exemplo-discriminador} ilustra a construção de discriminadores para o crime \textit{armas\_municoes}. O vocabulário ocupa posições estáveis na retina textual; o documento ativa as posições correspondentes às palavras observadas; e cada discriminador é uma combinação de posições que fornece evidência para a máscara do crime. O exemplo é didático e apresenta apenas dez posições; ao final da execução, a memória WNN continha 358 posições lexicais versionadas. Os cartões perfurados simplificados da figura mostram que cada crime ocupa um padrão distinto sobre as mesmas posições; eles não representam a quantidade total de bits nem os padrões completos da memória.

\begin{figure}[ht]
\centering
\begin{tikzpicture}[font=\scriptsize, >=latex]
  \node[draw, rounded corners, align=center, text width=.88\columnwidth, minimum height=.55cm, fill=gray!12] (text) {Trecho da notícia: ``apreensão de armas, munições, fuzil e carregadores''};
  \node[below=.34cm of text, align=center] (vocab) {\textbf{Vocabulário versionado}\\[-1mm]
  $a_1$ apreensão \quad $a_2$ armas \quad $a_3$ munições \quad $a_4$ arsenal\\[-1mm]
  $a_5$ fuzil \quad $a_6$ pistola \quad $a_7$ carregadores \quad $a_8$ lavagem\\[-1mm]
  $a_9$ dinheiro \quad $a_{10}$ ocultação};
  \draw[->, thick] (text.south) -- (vocab.north);
  \node[draw, rounded corners, below=.36cm of vocab, align=center, fill=blue!7, text width=.88\columnwidth] (vector) {\textbf{Vetor binário da notícia}\\[1mm]
  $\mathbf{x}(d)=[1,1,1,0,1,0,1,0,0,0]$};
  \draw[->, thick] (vocab.south) -- (vector.north);
  \node[draw, rounded corners, below=.36cm of vector, align=center, fill=green!9, text width=.88\columnwidth] (disc) {\textbf{Discriminadores da máscara \textit{armas\_municoes}}\\[.5mm]
  $d_1=\{a_1,a_2,a_3\}$ é ativado integralmente; $d_2$ é parcial e $d_3$ permanece inativo.};
  \draw[->, thick] (vector.south) -- (disc.north);
  \coordinate (c1) at ($(disc.south)+(-3.55cm,-.42cm)$);
  \coordinate (c2) at ($(disc.south)+(-.72cm,-.42cm)$);
  \coordinate (c3) at ($(disc.south)+(2.11cm,-.42cm)$);
  \foreach \card/\name in {c1/{armas e munições},c2/{contrabando},c3/{corrupção}} {
    \fill[black!88, rounded corners=1pt] ($(\card)+(-.08cm,.12cm)$) rectangle ($(\card)+(1.38cm,-.86cm)$);
    \node[font=\tiny, text=white] at ($(\card)+(.65cm,-.03cm)$) {\name};
    \foreach \x in {0,...,4} {
      \foreach \y in {0,...,2} {
        \fill[blue!25] ($(\card)+(.08cm+\x*.24cm,-.20cm-\y*.19cm)$) rectangle ++(.13cm,-.11cm);
      }
    }
  }
  \foreach \x/\y in {0/0,2/0,4/0,1/1,3/2} \fill[green!65] ($(c1)+(.08cm+\x*.24cm,-.20cm-\y*.19cm)$) rectangle ++(.13cm,-.11cm);
  \foreach \x/\y in {1/0,2/1,4/1,0/2,3/2} \fill[green!65] ($(c2)+(.08cm+\x*.24cm,-.20cm-\y*.19cm)$) rectangle ++(.13cm,-.11cm);
  \foreach \x/\y in {0/0,3/0,1/1,2/1,4/2} \fill[green!65] ($(c3)+(.08cm+\x*.24cm,-.20cm-\y*.19cm)$) rectangle ++(.13cm,-.11cm);
  \node[font=\tiny, align=center] at ($(disc.south)+(0,-1.52cm)$) {Cartões perfurados simplificados: verde = posição lexical associada à máscara do crime.};
\end{tikzpicture}
\caption{Exemplo didático de vetor binário, ativação de discriminadores e cartões perfurados simplificados das máscaras de crime.}
\label{fig:exemplo-discriminador}
\end{figure}

Uma coincidência parcial não é suficiente para considerar um discriminador integralmente ativado. No exemplo, a ocorrência de \textit{fuzil} e \textit{carregadores} ativa apenas parte das posições de $d_2$, pois o trecho não contém \textit{arsenal} nem \textit{pistola}; por isso, $d_2$ não produz a mesma evidência de $d_1$. Analogamente, como nenhuma palavra de $d_3$ ocorre no documento, esse discriminador permanece inativo. A decisão da memória considera a cobertura de cada discriminador e da máscara do crime, além da confiança e da margem entre classes, evitando que uma palavra isolada seja tomada como evidência conclusiva.

\subsection{Fundação temática: 10\% da base}

A fundação temática executa duas funções. A Função~1 consolida temas canônicos a partir dos clusters exploratórios. Inicialmente, os documentos são pré-processados, vetorizados e agrupados por HDBSCAN; em seguida, a similaridade do cosseno auxilia a aproximar clusters semanticamente equivalentes. A função interpreta os agrupamentos, controla atributos acidentais --- como localidades, órgãos e nomes de operações --- e define crimes canônicos que representam o fenômeno substantivo predominante.

Na execução, o pré-processamento linguístico foi aplicado aos 8.310 documentos antes da fundação e da classificação incremental. A representação preservou substantivos, verbos, adjetivos e expressões de domínio, removendo 497.688 de 1.240.262 tokens originais (redução de 40,13\%) e mantendo 742.574 tokens para as etapas analíticas.

A Função~2 gera discriminadores lexicais de crime. Para cada crime canônico, são selecionados marcadores substantivos e combinações de marcadores que compõem sua máscara vetorial WNN. A saída é um banco de discriminadores auditável, que registra o crime associado, os termos, as posições lexicais, a origem e o número de confirmações.

\subsection{Funções e instruções operacionais}

As funções do fluxo são definidas por contratos de entrada, instruções de operação e saídas estruturadas. Não se reproduzem os \textit{prompts} integralmente, pois eles incluem dados da rodada, mas a Tabela~\ref{tab:funcoes-instrucoes} registra as instruções que delimitam as decisões e permitem reproduzir o procedimento.

\begin{table}[ht]
\centering
\caption{Funções, instruções operacionais e saídas auditáveis.}
\label{tab:funcoes-instrucoes}
\begin{tabular}{>{\hyphenpenalty=10000\exhyphenpenalty=10000\emergencystretch=1em\arraybackslash}p{0.25\linewidth}>{\hyphenpenalty=10000\exhyphenpenalty=10000\emergencystretch=1em\arraybackslash}p{0.66\linewidth}}
\hline
\textbf{Função} & \textbf{Instrução operacional e saída} \\
\hline
1. Consolidação temática & Recebe resumos dos clusters, termos de domínio e evidências semânticas. Deve agrupar folhas equivalentes em crimes canônicos, ignorar localidades, órgãos, telefones, siglas e nomes de operações como possíveis classes. Pode aceitar, descartar ou colocar um cluster em quarentena; retorna rótulo canônico, clusters incluídos, termos de evidência, confiança e decisão. \\
2. Geração de discriminadores & Recebe os crimes aceitos, suas evidências e exemplos da fundação. Deve produzir combinações não ordenadas de termos substantivos, filtrando sinais acidentais e palavras genéricas. Retorna identificador, crime associado, termos, posições lexicais, peso, origem e justificativa de cada discriminador no banco de memória. \\
3. Classificação pela WNN & Recebe o corpo pré-processado, o título como apoio secundário, tags apenas como pistas e a versão corrente da memória. Deve projetar o documento no vetor binário, calcular cobertura, pontuação, confiança e margem das máscaras. Aceita um crime somente se os limiares configurados forem atendidos; caso contrário, encaminha o documento como residual e preserva as evidências ativadas. \\
4. Revisão residual por LLM & Recebe corpo, título, tags, classes canônicas permitidas, evidências WNN e sugestões por cosseno. Deve priorizar o sentido substantivo do corpo, usar o título como indício e não aceitar tags isoladamente. Confirma uma classe existente, propõe tema candidato, registra notícia rara ou indica quarentena para texto insuficiente. Retorna JSON com decisão, rótulo, confiança, evidência textual, justificativa, tema principal e marcadores secundários. \\
\hline
\end{tabular}
\end{table}

Após a revisão residual, um candidato não é promovido automaticamente. A governança compara suas evidências com as máscaras e os crimes já existentes; quando houver equivalência semântica, o caso é absorvido pela classe canônica correspondente. Por exemplo, ocorrências de pornografia infantil ou infantojuvenil são incorporadas a \textit{crimes\_contra\_criancas}, e não criam uma nova categoria. Somente candidatos recorrentes, textualmente distintos e confirmados podem ampliar a taxonomia.

\subsection{Classificação incremental: 90\% da base}

A Função~3 processa a reserva incremental em lotes. O corpo pré-processado de cada notícia é projetado pela retina textual no vetor $\mathbf{x}(d)$ e comparado às máscaras de crime. De forma abstrata, a pontuação de um crime pode ser expressa por

\begin{equation}
S(c,d)=\sum_{i=1}^{n_t}m_{ci}^{(t)}x_i(d).
\end{equation}

Na implementação, a pontuação considera as posições efetivamente ativadas, a cobertura da máscara, a confiança da decisão e sua margem em relação às alternativas. A notícia é aceita pela WNN somente quando os limiares de evidência, confiança e margem são atendidos. A decisão aceita registra o crime canônico, os marcadores ativados, as posições binárias, a versão da memória e a origem das evidências.

O tratamento de decisões pouco separadas inspira-se conceitualmente no 
\textit{bleaching} da WiSARD~\cite{carneiro2017weightless}. Na técnica original, o limiar de resposta é ajustado para distinguir discriminadores empatados ou muito próximos. Nesta proposta, o mecanismo não implementa o \textit{bleaching} nas RAMs: quando a maior pontuação não apresenta margem ou confiança suficiente, a Função~4 encaminha a notícia à LLM para uma revisão contextual de desempate. Assim, a memória permanece responsável pelos casos com evidência estável, enquanto a LLM atua somente onde a separação entre máscaras não é conclusiva.

\subsection{Revisão residual e atualização controlada}

A Função~4 revisa exclusivamente documentos residuais, ambíguos ou semanticamente inéditos. A revisão contextual pode confirmar um crime canônico, registrar uma notícia rara ou propor um tema candidato. Ela não substitui a memória WNN para os documentos já aceitos com evidência suficiente.

O crescimento da memória é controlado por uma política de vocabulário imutável e expansível por acréscimo. Uma posição já existente mantém seu significado e seu índice durante toda a execução. Antes de criar uma posição, o sistema normaliza o termo e consulta o mapa $\textit{termo}\rightarrow a_i$. Se o termo já existir, sua posição é reutilizada; se a evidência corresponder apenas a uma combinação de termos existentes, a máscara do crime é atualizada sem expandir o vocabulário.

Uma palavra inédita sugerida pela revisão residual é registrada inicialmente como variante pendente. Ela só é admitida como nova posição $a_{n_t+1}$ após pelo menos duas confirmações recorrentes. Quando aprovada, a nova posição é acrescentada no final de $V^{(t)}$, sem reordenar as posições anteriores. Assim, vetores históricos permanecem comparáveis mediante preenchimento de zeros à direita. A criação de um novo crime adiciona uma nova máscara vetorial sobre o vocabulário existente; ela não exige, necessariamente, uma nova posição lexical.

\subsection{Avaliação e auditabilidade}

A avaliação operacional mede a cobertura da WNN, a taxa de encaminhamento residual, o consumo de tokens pela LLM, a quantidade de discriminadores confirmados, os candidatos a novos temas e a evolução do vocabulário. Como as tags podem ser usadas como pistas auxiliares, elas não devem constituir a única referência de avaliação. A qualidade da classificação deve ser examinada por uma partição textual sem uso de tags ou por uma amostra validada manualmente.

Cada execução preserva as máscaras vetoriais dos crimes, o vocabulário versionado, os vetores e posições ativadas, as decisões residuais e as alterações de taxonomia. Esses artefatos permitem reconstruir a justificativa de uma classificação e verificar quando uma posição ou máscara foi incorporada à memória.

\section{Resultados}

Esta seção apresenta os diagnósticos operacionais da rodada concluída. A unidade experimental foi a notícia pública da Polícia Federal. A base reuniu 8.310 documentos, dos quais os 831 (10\%) mais antigos compuseram a fundação temática e os 7.479 (90\%) posteriores formaram a reserva incremental. Temporalmente, a fundação compreendeu notícias de 29/11/2011 até 30/07/2020; a reserva iniciou em 30/07/2020 e foi processada sem reordenamento cronológico em 15 lotes de aproximadamente 500 notícias. A coincidência da data de corte ocorre porque há múltiplas notícias publicadas em 30/07/2020; a ordem original do manifesto define a separação entre fundação e reserva nessa data.

\subsection{Clusterização da fundação temporal}

Na fundação de 831 notícias, a representação TF--IDF produziu vocabulário de 12.285 termos e matriz esparsa com 94.357 valores não nulos (densidade de 0,00924). Como estratégia operacional de agrupamento, a execução utilizou \textit{MiniBatch K-Means} como fallback, produzindo 28 clusters brutos, sem clusters de ruído e com coeficiente de silhueta de 0,288877. A similaridade do cosseno consolidou três grupos de clusters semanticamente próximos: os clusters 2, 9 e 23 em \textit{corrupcao\_recursos\_publicos} (113 notícias); os clusters 3 e 15 em \textit{crime\_organizado} (75); e os clusters 10, 18 e 19 em \textit{crimes\_contra\_criancas} (57). Após essas fusões, a fundação passou a conter 23 clusters para interpretação temática.

\subsection{Evolução da taxonomia de crimes}

A fundação temática consolidou 17 crimes canônicos. Durante a análise incremental, o tema \textit{ameacas\_e\_terrorismo} foi promovido, resultando em 18 categorias criminais na taxonomia final. O candidato operacional \textit{rara\_combate\_disseminacao\_pornografia} foi absorvido por \textit{crimes\_contra\_criancas}, pois suas evidências mencionam pornografia infantil ou infantojuvenil. A Tabela~\ref{tab:taxonomia-crimes} apresenta as categorias por etapa. A categoria operacional \textit{noticias\_raras} não integra a tabela, pois registra exceções para auditoria e não constitui crime canônico.

\begin{table}[ht]
\centering
\caption{Crimes canônicos por etapa de consolidação.}
\label{tab:taxonomia-crimes}
\begin{tabular}{p{0.16\linewidth}p{0.72\linewidth}}
\hline
\textbf{Etapa} & \textbf{Crime canônico ou tema promovido} \\
\hline
\multicolumn{2}{l}{\textit{Fundação temática (10\% da base): 17 crimes}} \\
\hline
1 & armas\_municoes \\
2 & contrabando\_descaminho \\
3 & corrupcao\_desvio\_recursos\_publicos \\
4 & crime\_organizado \\
5 & crimes\_ambientais \\
6 & crimes\_ciberneticos \\
7 & crimes\_contra\_criancas \\
8 & crimes\_eleitorais \\
9 & crimes\_migratorios \\
10 & crimes\_previdenciarios \\
11 & crimes\_sistema\_financeiro \\
12 & fraudes\_auxilios\_beneficios \\
13 & lavagem\_dinheiro \\
14 & moeda\_falsa \\
15 & radiodifusao\_clandestina \\
16 & trabalho\_escravo \\
17 & trafico\_drogas \\
\hline
\multicolumn{2}{l}{\textit{Promovido na reserva incremental (90\% da base): 1 tema}} \\
\hline
18 & ameacas\_e\_terrorismo \\
\hline
\end{tabular}
\end{table}

\subsection{Cobertura da memória e encaminhamento residual}

O principal indicador é a cobertura da memória WNN: a proporção de documentos da reserva classificados pela camada determinística antes de qualquer revisão por LLM. A Tabela~\ref{tab:resultados-cobertura} sintetiza os indicadores acumulados. A WNN classificou diretamente 6.031 das 7.479 notícias da reserva, atingindo cobertura de 80,64\%. Os 1.448 documentos restantes foram registrados como residuais e demandaram revisão contextual por LLM.

\begin{table}[ht]
\centering
\caption{Diagnósticos acumulados da rodada com fundação de 10\%.}
\label{tab:resultados-cobertura}
\small
\begin{tabular}{p{0.62\linewidth}r}
\hline
\textbf{Indicador} & \textbf{Valor} \\
\hline
Notícias na base completa & 8.310 \\
Notícias na fundação temática & 831 \\
Notícias na reserva incremental & 7.479 \\
Período da fundação temporal & 29/11/2011--30/07/2020 \\
Início da reserva incremental & 30/07/2020 \\
Lotes processados & 15 \\
Notícias classificadas diretamente pela WNN & 6.031 \\
Cobertura acumulada da WNN & 80,64\% \\
Documentos residuais pós-WNN & 1.448 \\
Chamadas residuais à LLM & 1.448 \\
Taxa de chamada à LLM & 19,36\% \\
Tokens consumidos nas chamadas residuais & 2.615.125 \\
Modelo da revisão residual & \textit{gpt-4.1-mini} \\
Média de tokens por chamada residual & 1.806 \\
\hline
\end{tabular}
\end{table}

Esse resultado indica que a LLM não precisou ser aplicada à totalidade da reserva. Caso a revisão por LLM fosse usada como classificador principal, haveria até 7.479 chamadas potenciais; na rodada observada, foram necessárias 1.448 chamadas residuais. Portanto, 6.031 chamadas potenciais (80,64\%) foram evitadas pela camada determinística. O consumo de 2.615.125 tokens pelo \textit{gpt-4.1-mini}, equivalente a 1.806 tokens por chamada em média, torna o custo computacional da camada residual mensurável e reproduzível. A redução de chamadas é um indicador de eficiência operacional, e não uma medida de qualidade semântica por si só.

\subsection{Interpretação operacional e auditabilidade}

Os documentos aceitos pela WNN preservam o crime canônico atribuído, os discriminadores ativados, as posições lexicais correspondentes, a versão da memória e a origem da evidência. Nos casos de baixa confiança, margem insuficiente ou competição entre máscaras, a decisão é deslocada para a revisão contextual da Função~4. Desse modo, os resultados permitem identificar tanto a parcela absorvida por evidência lexical recorrente quanto a parcela que exige interpretação adicional.

O uso dos diagnósticos de cobertura deve ser interpretado com cautela. Cobertura WNN não equivale a acurácia, pois a rodada não possui uma amostra de referência integralmente validada por especialistas. As tags institucionais foram tratadas como pistas auxiliares e não como verdade de referência. Além disso, a comparação com execuções anteriores é apenas operacional, pois as rodadas diferem em base, fração amostral, fundação temática e parâmetros. Assim, os resultados sustentam a viabilidade do fluxo e a redução do acionamento de LLM, mas não substituem uma avaliação futura de precisão, revocação e medida-F1.

\section{Conclusão}

Este trabalho abordou o problema de organizar e classificar incrementalmente uma base textual institucional em crescimento, reduzindo o custo de inferência por LLM e preservando auditabilidade, adaptação a novos temas e estabilidade taxonômica. Para isso, foi proposta uma metodologia que combina fundação temática formada pelos documentos históricos mais antigos, máscaras vetoriais de crimes em uma memória WNN e revisão contextual de documentos residuais por LLM.

Na rodada concluída, com 10\% da base destinados à fundação temática, a memória WNN classificou diretamente 80,64\% da reserva incremental. A LLM ficou restrita aos 19,36\% de casos com evidência insuficiente, baixa margem ou ambiguidade entre máscaras. Essa separação permite que documentos recorrentes sejam tratados por uma camada determinística, enquanto a interpretação mais custosa é reservada aos casos nos quais ela é necessária.

A contribuição do trabalho não se limita à classificação de uma coleção específica de notícias. A proposta organiza um ciclo auditável que transforma agrupamentos exploratórios em crimes canônicos, associa cada crime a uma máscara vetorial de evidências, registra as decisões da memória e encaminha incertezas para revisão residual. O vocabulário versionado e a política de acréscimo controlado preservam a comparabilidade das posições lexicais ao longo das atualizações.

Os resultados devem ser interpretados como diagnósticos operacionais internos. A qualidade da fundação temática, a representatividade temporal da amostra e a especificidade dos discriminadores influenciam a classificação. Sem anotação humana independente, a cobertura observada não permite afirmar acurácia. Como trabalhos futuros, recomenda-se validar uma amostra com especialistas, calcular precisão, revocação e medida-F1, comparar a proposta com classificadores supervisionados e LLM aplicada a toda a base e testar o procedimento em outras coleções institucionais.


\bibliographystyle{IEEEtran}
\bibliography{./trabalho_wnn/referencias}

\end{document}
